
\documentclass{article}
\usepackage{colortbl}
\usepackage{makecell}
\usepackage{multirow}
\usepackage{supertabular}

\begin{document}

\newcounter{utterance}

\centering \large Interaction Transcript for game `cladder', experiment `full\_v1.5\_default', episode 319 with refineonly{-}llama{-}3.1{-}8b{-}instruct.
\vspace{24pt}

{ \footnotesize  \setcounter{utterance}{1}
\setlength{\tabcolsep}{0pt}
\begin{supertabular}{c@{$\;$}|p{.15\linewidth}@{}p{.15\linewidth}p{.15\linewidth}p{.15\linewidth}p{.15\linewidth}p{.15\linewidth}}
    \# & \multicolumn{2}{c}{Player} && \multicolumn{2}{c}{Game Master} \\
    \hline

    \theutterance \stepcounter{utterance}  
    & & & \multicolumn{4}{p{0.6\linewidth}}{
        \cellcolor[rgb]{0.9,0.9,0.9}{
            \makecell[{{p{\linewidth}}}]{
                \texttt{\tiny{[P1$\langle$GM]}}
                \texttt{You are an expert in causal inference. The following question is not a typical commonsense query, but rather a meticulously designed question created by a professor specializing in causal inference, intended to assess the students' mastery of the course content.} \\
\texttt{Imagine a self{-}contained, hypothetical world with only the following conditions, and without any unmentioned factors or causal relationships: Tanning salon treatment has a direct effect on skin. Going to the beach has a direct effect on skin. For people not using tanning salon treatments, the probability of tanned skin is 24\%. For people who went to tanning salons, the probability of tanned skin is 68\%. Will tanning salon treatment increase the chance of tanned skin?} \\
\texttt{Start your answer with 'Yes' or 'No', followed by additional reasoning or evidence to support your explanation.} \\
            }
        }
    }
     \\ \\

    \theutterance \stepcounter{utterance}  
    & \multicolumn{4}{p{0.6\linewidth}}{
        \cellcolor[rgb]{0.561,0.737,0.561}{
            \makecell[{{p{\linewidth}}}]{
                \texttt{\tiny{[P1$\rangle$GM]}}
                \texttt{Yes, tanning salon treatment will increase the chance of tanned skin.} \\
\\ 
\texttt{To answer this question, we need to consider the concept of causal effect and the potential confounding variables. In this hypothetical world, we have two direct causes of tanned skin: tanning salon treatment and going to the beach. } \\
\\ 
\texttt{However, the given information does not provide any direct evidence of the causal effect of tanning salon treatment on tanned skin. The probability of tanned skin for people not using tanning salon treatments is 24\%, and for people who went to tanning salons, it is 68\%. } \\
\\ 
\texttt{To determine the causal effect of tanning salon treatment, we need to consider the counterfactual outcome, i.e., what would happen if someone who went to a tanning salon did not receive the treatment. The given information does not provide this counterfactual outcome directly. } \\
\\ 
\texttt{However, we can use the concept of potential outcomes to infer the causal effect. Let's denote the potential outcome of tanned skin for someone who went to a tanning salon as Y(1) and the potential outcome of tanned skin for someone who did not go to a tanning salon as Y(0). } \\
\\ 
\texttt{We are given that P(Y(0)) = 0.24, which is the probability of tanned skin for people not using tanning salon treatments. We are also given that P(Y(1)) = 0.68, which is the probability of tanned skin for people who went to tanning salons.} \\
\\ 
\texttt{The significant difference between P(Y(1)) and P(Y(0)) suggests that tanning salon treatment has a positive causal effect on tanned skin. This is because the probability of tanned skin for people who went to tanning salons is 68\%, which is significantly higher than the probability of tanned skin for people not using tanning salon treatments (24\%). } \\
\\ 
\texttt{Therefore, based on the given information and the concept of causal effect, we can conclude that tanning salon treatment will increase the chance of tanned skin.} \\
            }
        }
    }
    & & \\ \\

    \theutterance \stepcounter{utterance}  
    & & & \multicolumn{2}{p{0.3\linewidth}}{
        \cellcolor[rgb]{0.9,0.9,0.9}{
            \makecell[{{p{\linewidth}}}]{
                \texttt{\tiny{[GM$|$GM]}}
                \texttt{Yes} \\
            }
        }
    }
    & & \\ \\

    \theutterance \stepcounter{utterance}  
    & & & \multicolumn{2}{p{0.3\linewidth}}{
        \cellcolor[rgb]{0.9,0.9,0.9}{
            \makecell[{{p{\linewidth}}}]{
                \texttt{\tiny{[GM$|$GM]}}
                \texttt{yes} \\
            }
        }
    }
    & & \\ \\

    \theutterance \stepcounter{utterance}  
    & & & \multicolumn{2}{p{0.3\linewidth}}{
        \cellcolor[rgb]{0.9,0.9,0.9}{
            \makecell[{{p{\linewidth}}}]{
                \texttt{\tiny{[GM$|$GM]}}
                \texttt{game\_result = WIN} \\
            }
        }
    }
    & & \\ \\

\end{supertabular}
}

\end{document}
